%======================================================================
\chapter{Function spaces, the Poincar\'e inequality and C\'ea's lemma}\label{app:spaces}
%======================================================================
% Written for the author's note: "need of an introduction of functional
% spaces with examples and counterexamples; small explanation when
% Poincare appears; similar for Cea's theorem".

\framepart{Outline}
This appendix gives the minimum of functional analysis used in the book, in
the hands-on style of Strang~\cite{StrangFix1973,StrangCSE2007}: every
definition comes with its finite-dimensional analogue, an example and a
counterexample. Section~\ref{sec:sp-spaces} defines the spaces $L^2$, $H^1$
and $H^1_0$ and explains why they must be complete.
Section~\ref{sec:sp-traces} treats boundary values and loads.
Section~\ref{sec:sp-poincare} proves the Poincar\'e inequality and
Section~\ref{sec:sp-cea} C\'ea's lemma. For a complete treatment see
\cite{BrezisSp,EvansSp,AdamsSp}.

%======================================================================
\section{Spaces of finite energy}\label{sec:sp-spaces}
\index{function space}%
\paragraph{Norms and inner products.} On $\R^n$ the Euclidean inner product
$\vu\cdot\vv=\sum u_iv_i$ gives the norm $|\vu|$. For functions on a domain
$\Omega$, sums become integrals:
\begin{align*}
  (u,v)_{L^2}&=\int_\Omega uv\dd x, &
  \norm{v}_{L^2}^2&=\int_\Omega v^2\dd x,\\
  (u,v)_{H^1}&=\int_\Omega\bigl(uv+\nabla u\cdot\nabla v\bigr)\dd x, &
  \norm{v}_{H^1}^2&=\norm{v}_{L^2}^2+\norm{\nabla v}_{L^2}^2 .
\end{align*}
$L^2(\Omega)$ and $H^1(\Omega)$ are the functions for which these norms are
finite; $H^1_0(\Omega)$ is the subspace of $H^1$ of functions vanishing on
$\partial\Omega$. Derivatives are taken in the weak sense: $g=v'$ if
$\displaystyle \int g\varphi=-\int v\varphi'$ for all smooth $\varphi$ vanishing at the ends,
which is integration by parts used as a definition. A hat function has the
weak derivative $\pm1/h$, piecewise constant; a step function has none in
$L^2$ (its derivative is a Dirac mass).

\paragraph{Why completeness matters.}
\index{completeness}\index{Hilbert space}%
The direct method (Section~\ref{sec:var-direct}) and Lax--Milgram
(Section~\ref{sec:var-lm}) need limits of sequences to stay in the space.
$L^2$, $H^1$ and $H^1_0$ are complete: they are \emph{Hilbert spaces}.
Spaces of smooth functions are not. The functions
$v_n(x)=\sqrt{x^2+1/n^2}$ are smooth on $(-1,1)$ and converge in the $H^1$ norm
to $|x|$, which is not differentiable at $0$. A minimizer searched among
smooth functions may therefore not exist, while it exists in $H^1$; this is
why the weak form is posed in $H^1$ and not in $C^1$ or $C^2$.

\paragraph{Continuity depends on the dimension.}
\index{Sobolev embedding}%
In one dimension an $H^1$ function is continuous, and even
H\"older-continuous:
\[
  |v(x)-v(y)|=\Bigl|\int_y^xv'\Bigr|\le|x-y|^{1/2}\,\norm{v'}_{L^2}
\]
by Cauchy--Schwarz. In two and three dimensions this fails: $\ln|\ln r|$ is in
$H^1$ near the origin of $\R^2$ but unbounded (Exercise~\ref{exo:sp-loglog}).
Point values, and hence point loads and point constraints, make sense in 1D
(springs, bars, beams) but not in 2D or 3D, where a point force produces a
displacement with infinite energy (Exercise~\ref{exo:sp-point}).

%======================================================================
\section{Boundary values and loads}\label{sec:sp-traces}
\paragraph{Traces.}
\index{trace}%
An $L^2$ function has no boundary value: changing it on the boundary, a set of
zero volume, does not change it. An $H^1$ function has one, its \emph{trace},
and
\[
  \norm{v}_{L^2(\partial\Omega)}\le C_T\,\norm{v}_{H^1(\Omega)} .
\]
The traces of $H^1(\Omega)$ functions form the space $H^{1/2}(\partial\Omega)$,
``half a derivative'' smoother than $L^2$. Prescribed potentials or
displacements must therefore belong to $H^{1/2}$: a potential that jumps
along the boundary (two electrodes in contact) is not admissible, and the
corresponding energy is infinite.

\paragraph{Loads and the dual space.}
\index{dual space}%
A load is a linear form $L(v)$ on the space of test functions: the work of
the forces in the virtual displacement $v$. It must be continuous,
$|L(v)|\le\norm{L}\,\norm{v}$, i.e.\ belong to the dual space; for
$H^1_0(\Omega)$ the dual is denoted $H^{-1}(\Omega)$. Body forces in $L^2$ and
surface tractions in $L^2(\partial\Omega)$, or more generally in
$H^{-1/2}(\partial\Omega)$, are admissible. A point force $L(v)=Fv(\vx_0)$ is
admissible in 1D, where point values are continuous on $H^1$, and not in 2D
or 3D. The discrete analogue is a vector $\vf$ acting on $\vu$ by $\vf\T\vu$.

%======================================================================
\section{The Poincar\'e inequality}\label{sec:sp-poincare}
\index{Poincar\'e inequality}%
\begin{theorem}[Poincar\'e]\label{thm:poincare}
Let $\Omega$ be bounded. There is a constant $C_P$, depending only on $\Omega$,
such that
\[
  \norm{v}_{L^2(\Omega)}\le C_P\,\norm{\nabla v}_{L^2(\Omega)}\qquad\forall v\in H^1_0(\Omega).
\]
On an interval of length $L$ the best constant is $C_P=L/\pi$.
\end{theorem}
\begin{proof}[Proof on an interval]
Let $v\in H^1_0(0,L)$. Since $v(0)=0$,
$\displaystyle v(x)=\int_0^xv'(t)\dd t$, and by Cauchy--Schwarz
$\displaystyle v(x)^2\le x\int_0^Lv'^2$. Integrating over $(0,L)$ gives
$\displaystyle \int_0^Lv^2\le\frac{L^2}{2}\int_0^Lv'^2$, i.e.\ the inequality with
$C_P=L/\sqrt2$. For the best constant expand $v$ in the sine series
$v=\sum_kb_k\sin(k\pi x/L)$, which satisfies the boundary conditions:
$\displaystyle \int v^2=\frac L2\sum b_k^2$ and
$\displaystyle \int v'^2=\frac L2\sum(k\pi/L)^2b_k^2\ge(\pi/L)^2\int v^2$, with equality
for $v=\sin(\pi x/L)$.
\end{proof}

\paragraph{What it means.} A function pinned to zero on the boundary cannot be
large without slopes. The best constant is the inverse square root of the
smallest eigenvalue of $-\Delta$ with Dirichlet conditions,
$\displaystyle 1/C_P^2=\lambda_1=\min_v\int|\nabla v|^2/\int v^2$, the Rayleigh quotient of
\eqref{eq:rayleigh}. Physically $\lambda_1$ fixes the fundamental frequency of
a membrane or a string, and the Euler load of a column
(Section~\ref{sec:var-lm}).

\paragraph{Counterexamples.} Without a boundary condition the inequality fails:
constants have zero gradient. On $H^1(\Omega)$ it holds only for functions of
zero mean (Poincar\'e--Wirtinger inequality), which is why the pure Neumann
problem is solved up to a constant (Exercise~\ref{exo:neumann}). In matrix
terms, $\mat{B}\T\mat{B}$ is singular when the springs are not attached to a
wall: the rigid translation $\vu=(1,1,1)\T$ costs no energy.

%======================================================================
\section{C\'ea's lemma}\label{sec:sp-cea}
\index{C\'ea's lemma@C\'ea's lemma}\index{Galerkin method!orthogonality}%
\begin{theorem}[C\'ea]\label{thm:cea}
Let $a$ be bounded ($M$) and coercive ($\alpha$) on $V$, $u\in V$ the solution of
$a(u,v)=L(v)$ for all $v\in V$, and $u_h\in V_h\subset V$ the Galerkin solution,
$a(u_h,v_h)=L(v_h)$ for all $v_h\in V_h$. Then
\[
  \norm{u-u_h}\le\frac M\alpha\,\inf_{v_h\in V_h}\norm{u-v_h} .
\]
\end{theorem}
\begin{proof}
Subtracting the two equations for $v=v_h\in V_h$ gives the \emph{Galerkin
orthogonality} $a(u-u_h,v_h)=0$ for all $v_h\in V_h$. For any $v_h$, since
$u_h-v_h\in V_h$,
\[
  \alpha\norm{u-u_h}^2\le a(u-u_h,u-u_h)=a(u-u_h,u-v_h)\le M\norm{u-u_h}\norm{u-v_h}.
\]
Divide by $\norm{u-u_h}$ and take the infimum over $v_h$.
\end{proof}

\paragraph{What it means.} The Galerkin method makes no error of its own: up to
the factor $M/\alpha$ it delivers the best approximation that the space $V_h$
can offer. If $a$ is symmetric, Galerkin orthogonality says that $u_h$ is the
orthogonal projection of $u$ on $V_h$ for the energy inner product $a$; by
Pythagoras, $a(u-v_h,u-v_h)=a(u-u_h,u-u_h)+a(u_h-v_h,u_h-v_h)$, so $u_h$ is the
best approximation in the energy norm and the factor $M/\alpha$ disappears.
In matrix terms, the Galerkin solution solves the reduced system
$\mat{W}\T\mat{K}\mat{W}\,\vect{c}=\mat{W}\T\vf$, where the columns of
$\mat{W}$ span the trial space.

\paragraph{Convergence rate.} The infimum is bounded by choosing for $v_h$ the
interpolant $I_hu$ of $u$ with the hat functions. On each element of size
$h$, $|(u-I_hu)'|$ is controlled by $h\,|u''|$, so
$\norm{u-u_h}_{H^1}\le C\,h\,\norm{u''}_{L^2}$: linear elements converge at
order $h$ in energy, and at order $h^2$ in $L^2$ (Exercise~\ref{exo:conv}).

%======================================================================
\framepart{Summary}
\begin{gbox*}
\begin{itemize}
  \item \textbf{$L^2$, $H^1$, $H^1_0$} are the functions with finite mean
  square, finite energy, and finite energy with zero boundary values. They
  are complete (Hilbert spaces), which is what existence proofs need.
  \item \textbf{Examples and counterexamples:} hat functions are in $H^1$,
  steps are not; in 2D and 3D, $H^1$ functions may be unbounded and point
  forces have infinite energy.
  \item \textbf{Traces} of $H^1$ functions lie in $H^{1/2}(\partial\Omega)$;
  loads are elements of the dual space.
  \item \textbf{Poincar\'e:} on $H^1_0$ the gradient controls the function,
  $C_P=1/\sqrt{\lambda_1}$; the discrete analogue is the absence of rigid
  motions.
  \item \textbf{C\'ea:} the Galerkin solution is the best approximation up to
  $M/\alpha$, and exactly the best in the energy norm for symmetric problems.
\end{itemize}
\end{gbox*}

%======================================================================
\framepart{Exercises}

\exercise{Powers of $x$}\label{exo:sp-powers}
For which $\alpha>0$ is $x^\alpha$ in $L^2(0,1)$? In $H^1(0,1)$? In
$H^2(0,1)$?
\begin{solution}
$\displaystyle \int_0^1x^{2\alpha}<\infty$ for all $\alpha>-\tfrac12$, so always in $L^2$;
$\displaystyle \int_0^1\alpha^2x^{2\alpha-2}<\infty$ iff $\alpha>\tfrac12$;
$\displaystyle \int_0^1x^{2\alpha-4}<\infty$ iff $\alpha>\tfrac32$ (or $\alpha=1$, where the
second derivative vanishes).
\end{solution}

\exercise{An unbounded function of finite energy}\label{exo:sp-loglog}
In the disc $r<\tfrac12$ of $\R^2$, show that $v=\ln|\ln r|$ belongs to $H^1$
but is unbounded.
\begin{solution}
$|\nabla v|=1/(r|\ln r|)$, so
$\displaystyle \int|\nabla v|^2=2\pi\int_0^{1/2}\frac{\dd r}{r\ln^2r}=\frac{2\pi}{\ln2}<\infty$
(substitute $s=\ln r$). The function itself is square integrable and tends to
$+\infty$ as $r\to0$.
\end{solution}

\exercise{A point force on a membrane}\label{exo:sp-point}
A membrane of tension $T$ occupies the disc $r<R$, is fixed on $r=R$ and
loaded by a point force $F$ at its centre. (a) Show that $\displaystyle u=\frac{F}{2\pi T}\ln(R/r)$
satisfies the equilibrium $-T\Delta u=0$ for $r>0$ and carries the total
force $F$. (b) Compute the energy in the annulus $\epsilon<r<R$ and let
$\epsilon\to0$. (c) Compare with a string loaded at one point.
\begin{solution}
(a) $\Delta\ln r=0$ for $r>0$, and the flux through a circle is
$-2\pi rT\,u'(r)=F$. (b) $\displaystyle \tfrac T2\int|\nabla u|^2=\frac{F^2}{4\pi T}\ln(R/\epsilon)\to\infty$:
infinite energy, and infinite displacement at the load point. (c) A string
loaded at one point takes a triangular shape of finite energy: in 1D point
values are continuous on $H^1$ and point forces are admissible. In 2D a
force must be spread over an area, however small.
\end{solution}

\exercise{The Poincar\'e constant by the Rayleigh quotient}\label{exo:sp-rayleigh}
Use the trial function $v=x(1-x)$ in the Rayleigh quotient
$\displaystyle \int_0^1v'^2/\int_0^1v^2$ to bound $\lambda_1$ on $(0,1)$ from above, and
compare with $\pi^2$.
\begin{solution}
$\displaystyle \int_0^1(1-2x)^2=\tfrac13$ and $\displaystyle \int_0^1x^2(1-x)^2=\tfrac1{30}$, so the
quotient is $10\ge\pi^2\approx9.87$: an error of $1.3\%$ with a single
parabola. This is the Rayleigh--Ritz method for eigenvalues, used later for
buckling loads.
\end{solution}

\exercise{Galerkin is exact at the nodes in 1D}\label{exo:sp-nodal}
For $-u''=f$ on $(0,1)$, $u(0)=u(1)=0$, show that the linear finite element
solution $u_h$ equals the exact solution at the nodes, so that $u_h=I_hu$
and the error is exactly the interpolation error of C\'ea's lemma.
\begin{solution}
Galerkin orthogonality reads $\displaystyle \int_0^1(u-u_h)'\varphi_i'=0$ for every hat
function. Since $\varphi_i'=\pm1/h$ on the two elements of node $i$, this is
$\displaystyle \frac1h\bigl[(e(x_i)-e(x_{i-1}))-(e(x_{i+1})-e(x_i))\bigr]=0$ for the error
$e=u-u_h$, with $e(0)=e(1)=0$: the nodal errors satisfy the discrete Laplace
equation with zero end values, hence vanish. This special property of 1D
explains the high nodal accuracy observed in Exercise~\ref{exo:conv}.
\end{solution}

%======================================================================
\begin{chapterbib}{9}
\bibitem{AdamsSp} R.~A.~Adams, J.~J.~F.~Fournier, \emph{Sobolev Spaces}, 2nd ed., Academic Press, 2003.
\bibitem{BrezisSp} H.~Brezis, \emph{Functional Analysis, Sobolev Spaces and Partial Differential Equations}, Springer, 2011.
\bibitem{EvansSp} L.~C.~Evans, \emph{Partial Differential Equations}, 2nd ed., AMS, 2010.
\bibitem{StrangCSE2007} G.~Strang, \emph{Computational Science and Engineering}, Wellesley-Cambridge Press, 2007.
\bibitem{StrangFix1973} G.~Strang, G.~J.~Fix, \emph{An Analysis of the Finite Element Method}, Prentice-Hall, 1973.
\end{chapterbib}
